
\documentclass[12pt]{article}

\usepackage{geometry}
\usepackage{titling}
\usepackage{amsmath,amssymb}
\usepackage{fontspec}
\usepackage{unicode-math}
\usepackage{array}
\usepackage{colortbl}
\usepackage[dvipsnames]{xcolor}
\usepackage[most]{tcolorbox}

\geometry{a4paper, top=2.5cm, bottom=2.5cm, left=2.6cm, right=2.5cm}
\setlength{\droptitle}{-3cm}

\setmonofont{DejaVu Sans Mono}
\setmainfont{Libertinus Serif}
\setsansfont{Libertinus Sans}
\setmathfont{STIXTwoMath-Regular.otf}

% Define the custom vector command used in the document
\newcommand{\MVec}[1]{\mathbf{#1}}

\title{Application of Euclidean Distance \vspace*{4pt}\\ in \vspace*{4pt}\\K-Nearest Neighbor Regression}
\author{Devi Prasad M}
\date{August 2026}

\begin{document}
\pagecolor{yellow!25!brown!30}

\maketitle

\section*{An Example: Finding the K-Nearest Neighbors}
Consider a two-dimensional feature space and the query point:
\[
\MVec{x}_q = (5,5).
\]
Suppose our training dataset contains the following $12$ points:
\[
\begin{array}{c|c|c}
\text{Point} & \MVec{x}_i & y_i \\ \hline
\MVec{x}_1 & (1,2) & 10 \\
\MVec{x}_2 & (2,8) & 16 \\
\MVec{x}_3 & (3,4) & 12 \\
\MVec{x}_4 & (4,6) & 14 \\
\MVec{x}_5 & (5,3) & 11 \\
\MVec{x}_6 & (6,5) & 15 \\
\MVec{x}_7 & (7,7) & 17 \\
\MVec{x}_8 & (8,4) & 18 \\
\MVec{x}_9 & (9,6) & 20 \\
\MVec{x}_{10} & (2,3) & 9 \\
\MVec{x}_{11} & (7,3) & 16 \\
\MVec{x}_{12} & (4,3) & 13
\end{array}
\]
We choose:
\[
K=3.
\]

\subsection*{Step 1: Measure the Distance to Every Point}
For the query $\MVec{x}_q=(5,5)$,
\[
\mathbf{dist}(\MVec{x}_q,\MVec{x}_i)
=
\sqrt{
(5-x_i^{(1)})^2+
(5-x_i^{(2)})^2
}.
\]
We calculate the distance to \textit{every} one of the $12$ training points:
\[
\begin{array}{c|c|c@{\qquad}c|c|c}
\text{Point} & y_i & \text{Distance}
& \text{Point} & y_i & \text{Distance} \\ \hline
\MVec{x}_1 & 10 & 5.00 & \MVec{x}_7 & 17 & 2.83 \\
\MVec{x}_2 & 16 & 4.24 & \MVec{x}_8 & 18 & 3.16 \\
\MVec{x}_3 & 12 & 2.24 & \MVec{x}_9 & 20 & 4.12 \\
\MVec{x}_4 & 14 & 1.41 & \MVec{x}_{10} & 9 & 3.61 \\
\MVec{x}_5 & 11 & 2.00 & \MVec{x}_{11} & 16 & 2.83 \\
\MVec{x}_6 & 15 & 1.00 & \MVec{x}_{12} & 13 & 2.24
\end{array}
\]
\[
\boxed{
\text{Measure: Every training point is compared with the query.}
}
\]

\subsection*{Step 2: Select the $K$ Nearest Neighbors}
Since $K=3$, select the three training points having the smallest Euclidean distances:
\[
\begin{array}{c|c|c}
\text{Neighbor} & \text{Target } y_i & \mathbf{dist}(\MVec{x}_q,\MVec{x}_i) \\ \hline
\MVec{x}_6=(6,5) & 15 & 1.00 \\
\MVec{x}_4=(4,6) & 14 & 1.41 \\
\MVec{x}_5=(5,3) & 11 & 2.00
\end{array}
\]
Therefore,
\[
N_3(\MVec{x}_q)
=
\left\{
(\MVec{x}_6,15),
(\MVec{x}_4,14),
(\MVec{x}_5,11)
\right\}.
\]

\subsection*{Step 3: Predict}
The KNN regression prediction is the arithmetic mean of the target values of the three nearest neighbors:
\[
\hat{y}(\MVec{x}_q)
=
\frac{15+14+11}{3}
=
\frac{40}{3}
\approx
13.33.
\]
Thus the KNN estimate for the unknown target at $\MVec{x}_q=(5,5)$ is:
\[
\boxed{
\hat{y}(\MVec{x}_q)\approx 13.33
}
\]

\subsection*{The Essential Idea}
The example illustrates the three stages of KNN regression:
\[
\boxed{
\begin{array}{c}
\text{Measure}\\
\text{Euclidean distance}
\end{array}
}
\quad
\longrightarrow
\quad
\boxed{
\begin{array}{c}
\text{Select}\\
K\text{ nearest neighbors}
\end{array}
}
\quad
\longrightarrow
\quad
\boxed{
\begin{array}{c}
\text{Predict}\\
\text{average their target values}
\end{array}
}
\]

Euclidean distance provides a geometric notion of proximity
based on the inner product. KNN uses this geometry to construct a
local estimate of an unknown function, based on the assumption
that nearby points tend to have similar target values.

\[
    \boxed{\text{Inner Product}}
    \
    \longrightarrow
    \
    \boxed{\text{Norm}}
    \
    \longrightarrow
    \
    \boxed{\text{Distance}}
    \
    \longrightarrow
    \
    \boxed{\text{Nearest Neighbor}}
    \
    \longrightarrow
    \
    \boxed{\text{Prediction}}
\]

\end{document}